\documentclass[preview]{standalone}

\usepackage{amsmath}
\usepackage{amssymb}
\usepackage{stellar}
\usepackage{definitions}

\begin{document}

\id{primary-rational-canonical-forms}
\genpage

\section{Primary canonical form}

\begin{snippettheorem}{endomorphism-primary-canonical-form-theorem}{Primary canonical form}
    Let \(V\) be a finite-dimensional \vectorspace over the \field \(K\),
    and let \(\varphi\colon V\fromto V\) be a
    \linearendomorphism. There exist monic
    \irreducibleelement[irreducible polynomials]
    \(f_1,\ldots,f_t\in K[X]\) and integers \(n_1,\ldots,n_t\geq1\)
    such that, in a suitable basis,
    \[
        [\varphi]
        =
        \begin{pmatrix}
            C(f_1^{n_1})&0&\cdots&0\\
            0&C(f_2^{n_2})&\cdots&0\\
            \vdots&\vdots&\ddots&\vdots\\
            0&0&\cdots&C(f_t^{n_t})
        \end{pmatrix}.
    \]
    The companion blocks \(C(f_i^{n_i})\) are uniquely determined up
    to reordering. This block matrix is a \emph{primary canonical form}
    of \(\varphi\).
\end{snippettheorem}

\begin{snippetproof}{endomorphism-primary-canonical-form-theorem-proof}{endomorphism-primary-canonical-form-theorem}{Primary canonical form}
    The \endomorphismmodule is finitely generated and torsion.
    Since \(K[X]\) is a \principalidealdomain, primary decomposition gives
    \[
        V\moduleisomorphic
        K[X]/(f_1^{n_1})
        \oplus\cdots\oplus
        K[X]/(f_t^{n_t}),
    \]
    where the \(f_i\) are monic irreducible polynomials. For each cyclic
    summand choose a generator \(v_i\), and put
    \(m_i=\deg(f_i^{n_i})\). The corresponding invariant subspace has
    basis
    \[
        \mathcal B_i
        =
        \{v_i,\varphi(v_i),\ldots,\varphi^{m_i-1}(v_i)\},
    \]
    in which the restriction of \(\varphi\) is represented by
    \(C(f_i^{n_i})\). Concatenating these bases produces the displayed
    block-diagonal matrix. Uniqueness follows from uniqueness of the
    elementary divisors of the \(K[X]\)-module.
\end{snippetproof}

\begin{snippetcorollary}{matrix-primary-canonical-form}{Primary canonical form for matrices}
    Every \(A\in\matrices_n(K)\) is similar to a block-diagonal matrix
    whose blocks are
    \[
        C(f_1^{n_1}),\ldots,C(f_t^{n_t}),
    \]
    where the \(f_i\in K[X]\) are monic irreducible polynomials and
    \(n_i\geq1\). The multiset of blocks is unique.
\end{snippetcorollary}

\begin{snippetproof}{matrix-primary-canonical-form-proof}{matrix-primary-canonical-form}{Primary canonical form for matrices}
    Associate with \(A\) the \linearendomorphism
    \[
        \varphi_A\colon K^n\fromto K^n,
        \qquad
        v\mapsto Av.
    \]
    Apply the primary canonical form theorem to \(\varphi_A\). Choosing
    a new basis replaces \(A\) by a similar matrix, and every similarity
    arises from such a change of basis.
\end{snippetproof}

\section{Rational canonical form}

\begin{snippettheorem}{endomorphism-rational-canonical-form-theorem}{Rational canonical form}
    Let \(V\) be a finite-dimensional \vectorspace over \(K\), and let
    \(\varphi\colon V\fromto V\) be a \linearendomorphism. There exist
    unique monic nonconstant polynomials
    \(d_1,\ldots,d_s\in K[X]\) such that
    \[
        d_1\divides d_2\divides\cdots\divides d_s
    \]
    and, in a suitable basis,
    \[
        [\varphi]
        =
        \begin{pmatrix}
            C(d_1)&0&\cdots&0\\
            0&C(d_2)&\cdots&0\\
            \vdots&\vdots&\ddots&\vdots\\
            0&0&\cdots&C(d_s)
        \end{pmatrix}.
    \]
    This matrix is the \emph{rational canonical form} of \(\varphi\);
    the \(d_i\) are its \invariantfactor[invariant factors].
\end{snippettheorem}

\begin{snippetproof}{endomorphism-rational-canonical-form-theorem-proof}{endomorphism-rational-canonical-form-theorem}{Rational canonical form}
    Invariant-factor decomposition of the associated torsion module
    gives
    \[
        V\moduleisomorphic
        K[X]/(d_1)\oplus\cdots\oplus K[X]/(d_s),
        \qquad
        d_1\divides\cdots\divides d_s.
    \]
    The monic invariant factors are unique. Choosing a generator of each
    summand gives an invariant cyclic subspace on which \(\varphi\) has
    matrix \(C(d_i)\). Concatenating the corresponding cyclic bases
    produces the rational canonical form.
\end{snippetproof}

\begin{snippetcorollary}{matrix-similarity-rational-canonical-form}{Similarity criterion}
    Two matrices \(A,B\in\matrices_n(K)\) are similar
    \ifandonlyif they have the same \rationalcanonicalform.
    Equivalently, they are similar \ifandonlyif their
    \primarycanonicalform[primary canonical forms] have the same
    companion blocks, up to reordering.
\end{snippetcorollary}

\begin{snippetproof}{matrix-similarity-rational-canonical-form-proof}{matrix-similarity-rational-canonical-form}{Similarity criterion}
    Similar matrices represent the same \linearendomorphism in different
    bases, so they have the same associated \(K[X]\)-module and hence
    the same invariant factors. Conversely, if both matrices have the
    same rational canonical form \(C\), each is similar to \(C\), and
    therefore they are similar to each other. The primary formulation
    follows in the same way from uniqueness of elementary divisors.
\end{snippetproof}

\section{Characteristic polynomial}

\begin{snippetproposition}{canonical-forms-characteristic-polynomial}{Characteristic polynomial from canonical blocks}
    If the primary canonical form of \(\varphi\) has blocks
    \[
        C(f_1^{n_1}),\ldots,C(f_t^{n_t}),
    \]
    then
    \[
        \chi_\varphi
        =
        f_1^{n_1}\cdots f_t^{n_t}.
    \]
    If its rational canonical form has blocks
    \(C(d_1),\ldots,C(d_s)\), then
    \[
        \chi_\varphi=d_1\cdots d_s.
    \]
\end{snippetproposition}

\begin{snippetproof}{canonical-forms-characteristic-polynomial-proof}{canonical-forms-characteristic-polynomial}{Characteristic polynomial from canonical blocks}
    The characteristic polynomial of a block-diagonal matrix is the
    product of those of its blocks, and the characteristic polynomial of
    \(C(f)\) is \(f\).
\end{snippetproof}

\end{document}
